\documentclass[pdflatex,sn-mathphys-num]{sn-jnl}
\usepackage{fix-cm}

\usepackage{graphicx}%
\usepackage{multirow}%
\usepackage{amsmath,amssymb,amsfonts}%
\usepackage{amsthm}%
\usepackage{mathrsfs}%
\usepackage[title]{appendix}%
\usepackage{xcolor}%
\usepackage{textcomp}%
\usepackage{manyfoot}%
\usepackage{booktabs}%
\usepackage{algorithm}%
\usepackage{algorithmicx}%
\usepackage{algpseudocode}%
\usepackage{listings}%
\setcounter{secnumdepth}{4}
\setcounter{tocdepth}{4}
\hypersetup{bookmarksdepth=4}

\theoremstyle{thmstyleone}%

\newtheorem{theorem}{Theorem}[section]
\newtheorem{lemma}[theorem]{Lemma}
\newtheorem{proposition}[theorem]{Proposition}
\newtheorem{corollary}[theorem]{Corollary}
\theoremstyle{definition}
\newtheorem{conjecture}[theorem]{Conjecture}
\newtheorem{observation}[theorem]{Empirical Observation}
\theoremstyle{plain}

\raggedbottom

\begin{document}

\title[A Symbolic Prefix Decomposition of the Collatz Map and Its Empirical Determination of Per-Class Statistics]{A Symbolic Prefix Decomposition of the Collatz Map and Its Empirical Determination of Per-Class Statistics}

\author[1]{\fnm{Nathan} \sur{Humphrey}}\email{mr.nathanhumphrey@gmail.com}
\affil[1]{\orgdiv{Founder}, \orgname{Resolve Research}, \orgaddress{\street{Kensington Way}, \city{Annapolis}, \postcode{21403}, \state{MD}, \country{USA}}}

\keywords{Collatz map, $3x+1$ problem, total stopping time, Syracuse parity vector, random-walk heuristic, experimental number theory}
\pacs[MSC Classification]{11B83, 37A45, 11Y55, 60G50}



\maketitle
\begin{abstract}
\noindent
Fix $k \geq 1$ and an odd residue $r \bmod 2^k$. Write $n = 2^k m + r$ and track the Collatz iteration symbolically as a linear state $a_t \cdot m + c_t$ initialized at $(a_0, c_0) = (2^k, r)$. While $a_t$ is even the parity of $a_t m + c_t$ depends only on $c_t \bmod 2$, so the iteration's behavior in this phase is determined by $r$ alone. We call this the \emph{prefix iteration} and study its terminal value $a_{\text{final}}$, reached when $a_t$ first becomes odd.

The contribution is in two parts. \emph{Algebraic.} The prefix iteration consumes exactly $k$ even steps and at most $k$ odd steps, terminating at $a_{\text{final}} \in \{3^j : 1 \leq j \leq k\}$ (Theorem~\ref{thm:termination}). The block decomposition of the prefix trace is the Syracuse parity vector of $r$ (Lemma~\ref{lem:syracuse-equivalence}), and the count of odd classes mod $2^k$ with $a_{\text{final}} = 3^j$ is $\binom{k-1}{j-1}$ via Terras's bijection (Lemma~\ref{lem:j-binomial}). Defining $\alpha_{\text{det}}(r) := \text{prefix\_steps}(r) + K_h \log(a_{\text{final}}(r)/2^k)$ with the random-walk constant $K_h = 3/\log(4/3)$, we prove $\langle \alpha_{\text{det}} \rangle = \log 6 / \log(4/3)$ exactly, independent of $k$ (Theorem~\ref{thm:alpha-det-closed-form}).

\emph{Experimental.} At modular resolutions $k \in \{6, 7, 8, 9\}$ and $N$ up to $2^{27}$, the per-class intercept $\alpha(r)$ of total stopping time $\sigma$ against $\log n$ is predicted by $\alpha_{\text{det}}(r)$ with regression $R^2 \geq 0.985$ and residual signal-to-noise ratio $\mathrm{SD}(\text{residual})/\mathrm{mean}(\mathrm{SE}) \in [0.91, 0.99]$. Per-class central moments of $\sigma$ residuals cluster by $a_{\text{final}}$ into exactly $k$ groups at every tested $k$. These regularities motivate Conjecture~\ref{conj:prefix-parameterization}: that the conditional distribution $\sigma \mid r \bmod 2^k$ is parameterized by $a_{\text{final}}$ alone at arbitrary $k$.

The prefix decomposition is a constructive sharpening of Lemma 4 of Terras: where Terras gives an asymptotic ratio for the $k$th iterate's residue-class structure, we retain the symbolic state through the deterministic prefix and expose $a_{\text{final}}$ as a closed-form covariate at finite $k$.
\end{abstract}

\noindent\textbf{Keywords:} Collatz map; $3x+1$ problem; total stopping time; Syracuse parity vector; random-walk heuristic; experimental number theory.

\noindent\textbf{2020 Mathematics Subject Classification:} 11B83 (primary); 37A45, 11Y55, 60G50 (secondary).

\section{Introduction}
For odd $n$ with $r = n \bmod 2^k$, the Collatz iteration starting from $n$ can be tracked symbolically. Writing $n = 2^k m + r$ and following the symbolic state $a_t \cdot m + c_t$ from $(a_0, c_0) = (2^k, r)$ under the Collatz map, the state's parity is determined by $r$ alone while $a_t$ remains even. We call this the \emph{prefix iteration} and write $a_{\text{final}}$ for the value of $a_t$ at the first step where it becomes odd. By Theorem~\ref{thm:termination}, $a_{\text{final}} \in \{3^j : 1 \leq j \leq k\}$, and the number of odd-step applications consumed is at most $k$.
 
The contribution has an algebraic part and an experimental part.
 
\subsection{Algebraic content (\S\ref{sec:prefix})} The principal results are: the termination theorem (Theorem~\ref{thm:termination}, the prefix iteration consumes exactly $k$ even steps and at most $k$ odd steps, terminating at $a_{\text{final}} \in \{3^j\}$); the binomial $j$-distribution lemma (Lemma~\ref{lem:j-binomial}, the count of odd classes mod $2^k$ with exactly $j$ odd-step applications is $\binom{k-1}{j-1}$, proven via the Syracuse parity-vector bijection of \citet{terras1976}); and the closed-form theorem $\langle \alpha_{\text{det}} \rangle = \log 6 / \log(4/3)$ (Theorem~\ref{thm:alpha-det-closed-form}, exact in the prefix algebra with the random-walk constant $K_h$ as the only heuristic input).
 
\subsection{Experimental content (\S\ref{sec:empirics})} At modular resolutions $k \in \{6, 7, 8, 9\}$ and $N$ up to $2^{27}$, two regularities hold. First, the per-class intercept $\alpha(r)$ of $\sigma$ against $\log n$ is predicted by $\alpha_{\text{det}}(r)$ with regression $R^2 \geq 0.985$. The residual signal-to-noise ratio $\mathrm{SD}(\text{residual})/\mathrm{mean}(\mathrm{SE})$ falls in $[0.91, 0.99]$ across the four tested resolutions. Per-class residuals scale with sampling noise alone. Second, per-class central moments of $\sigma$ residuals cluster by $a_{\text{final}}$, with exactly $k$ distinct clusters at every tested $k$. These findings validate the random-walk embedding $K_h$ underlying $\alpha_{\text{det}}$ on per-class $\sigma$ statistics at the tested resolutions. They motivate Conjecture~\ref{conj:prefix-parameterization}, that the conditional distribution $\sigma \mid r \bmod 2^k$ is parameterized by $a_{\text{final}}$ alone at arbitrary $k$.
 
The closed form $\langle \alpha_{\text{det}} \rangle = \log 6 / \log(4/3)$ is Theorem~\ref{thm:alpha-det-closed-form}, not an experimental claim; \S\ref{sec:closed-form-verification} reports the corresponding numerical sanity check.
 
\subsection{Relation to Terras (1976)} The prefix decomposition is a constructive sharpening of \citet{terras1976} Lemma 4. Terras's Lemma 4 establishes the asymptotic identity $S_k \approx S_0 \cdot 3^{d(k)} / 2^k$, where $S_k$ is the residue-class structure at the $k$th iterate and $d(k)$ is the number of odd-step applications among the first $k$ Collatz steps. The prefix decomposition retains the symbolic state $(a, c)$ through the deterministic prefix and terminates at $a_{\text{final}} = 3^j$ explicitly. Terras's lemma gives the asymptotic content; the prefix decomposition gives the constructive content at finite $k$. Section~\ref{sec:terras-relation} presents this relationship in full so that the contribution is bounded by what the decomposition adds to Terras's results: the constructive availability of $a_{\text{final}}$ and $\text{prefix\_steps}$ as closed-form predictors of per-class $\sigma$ statistics.
 
The independent parallel work of \citet{chang2026a, chang2026b} obtains a related decomposition of the burst-gap indicator at $K \in \{3, 4, 5\}$ via block total variation. Chang's framework addresses a different statistical question than the prefix decomposition introduced here.
 
\subsection{Organization} Section~\ref{sec:background} places the decomposition within the standard Collatz random-walk heuristic, \citet{terras1976} Lemma 4, and the symbolic approaches to Collatz residue dynamics. Section~\ref{sec:prefix} defines the prefix decomposition, states and proves the termination theorem (\S\ref{sec:termination}), establishes the realization of all $k$ values of $a_{\text{final}}$ as a corollary (\S\ref{sec:realization}), defines $\alpha_{\text{det}}$, and proves $\langle \alpha_{\text{det}} \rangle = \log 6 / \log(4/3)$ via the binomial $j$-distribution lemma (\S\ref{sec:closed-form}). Section~\ref{sec:empirics} presents the empirical regularities at $k \in \{6, 7, 8, 9\}$, $N \leq 2^{27}$ that validate the random-walk embedding underlying $\alpha_{\text{det}}$. Section~\ref{sec:terras-relation} situates the decomposition relative to \citet{terras1976}. Section~\ref{sec:discussion} states the conjecture the experimental regularities motivate, the analytical work that would strengthen it, and the $qx+1$ extension. Section~\ref{sec:conclusion} concludes.

\section{Background}
\label{sec:background}
\subsection{Collatz dynamics and the standard heuristic}
The Collatz map $T(n) = 3n + 1$ if $n$ is odd, $n/2$ if $n$ is even, generates the sequence $n \to T(n) \to T(T(n)) \to \cdots$. The total stopping time $\sigma(n)$ is the number of Collatz steps from $n$ to $1$. For random odd $n$, the standard random-walk heuristic predicts $\mathbb{E}[\sigma \mid n] \approx K_h \cdot \log(n)$ where $K_h = 3/ \log(4/3) \approx 10.4282$ \cite{lagarias1985}. The heuristic derives from a random-walk argument: the Syracuse compression $n \mapsto (3n + 1)/2^v$ has expected $\log$-multiplier $\log(3) - \langle v \rangle \cdot \log 2 = - \log(4/3)$ when $v$ is treated as $\text{Geometric}(\frac{1}{2})$, and the inverse of this drift gives $K_h$ \citep{lagarias2010, sinai2003}. Whether $\sigma$ has additional residue-class structure beyond the heuristic single-slope model is the question this paper addresses for $\sigma$ on odd integers.

\subsection{Terras 1976 Lemma 4}
\citet{terras1976} established the asymptotic identity for the $k$th Collatz iterate's residue-class structure: $$S_k \approx S_0 \cdot 3^{d(k)}/2^k$$ where $d(k)$ is the number of odd-step applications consumed among the first $k$ Collatz steps starting from the residue class. Terras's lemma gives the asymptotic ratio between successive iterates' residue-class structure. It does not retain the symbolic state $(a, c)$ past the asymptotic identity. The prefix decomposition introduced in \S3 is a sharpening of Terras Lemma 4: we retain the full symbolic state through the deterministic prefix, terminate at $a_{\text{final}} = 3^j$ explicitly, and use this terminal value as a closed-form covariate for per-class $\sigma$ statistics rather than only for the asymptotic step ratio.

\subsection{Symbolic and algebraic approaches to Collatz residue dynamics}
Symbolic and algebraic approaches to Collatz residue dynamics have a substantial history. The natural setting is the $(\mathbb{Z}/3^k\mathbb{Z})^*$ tower, with reduction maps connecting modular resolutions, and the 2-adic structure governing the parity dynamics of even-step compressions. A comprehensive bibliography of this literature is collected in \citet{lagarias2010}. The prefix decomposition introduced here is restricted to the deterministic prefix region, the part of the Collatz iteration where the symbolic state's parity is determined by $r$ alone, and does not address the post-prefix dynamics, which are stochastic from the per-class perspective.

\section{The prefix decomposition: definition and algebraic properties}
\label{sec:prefix}
\subsection{The symbolic state and parity rules}
\label{subsec:prefix-recall}
Fix an odd integer $n$ with residue $r = n \bmod 2^k$, and write $n = 2^k \cdot m + r$, where $m \geq 0$ is an integer (with $m = 0$ corresponding to $n = r$, the smallest representative of the class). Define the symbolic state at step $t$ of the Collatz iteration as $a_t \cdot m + c_t$, with initial state $(a_0, c_0) = (2^k, r)$. The state is a linear function of $m$, with coefficients $(a_t, c_t)$ determined by the iteration's history.

The Collatz iteration's parity at step $t$ is determined by $a_t \cdot m + c_t \bmod 2$. While $a_t$ is even, the parity is $c_t \bmod 2$, independent of $m$. While $a_t$ is odd, the parity depends on $m$, and the iteration's behavior is no longer determined by $r$ alone.

The state transitions are:
\begin{enumerate}[(i)]
    \item If $a_t$ is even and $c_t$ is even: $a_t \cdot m + c_t$ is even for all $m$. The iteration applies $n \rightarrow n/2$, giving $(a_{t+1}, c_{t+1}) = (a_t/2, c_t/ 2)$.
    \item If $a_t$ is even and $c_t$ is odd: $a_t \cdot m + c_t$ is odd for all $m$. The iteration applies $n \rightarrow 3n+1$, giving $(a_{t+1}, c_{t+1}) = (3 \cdot a_t, 3 \cdot c_t + 1)$.
    \item If $a_t$ is odd: parity depends on $m$. The prefix iteration terminates.
\end{enumerate}

We define the prefix iteration as the sequence of state transitions (i) and (ii) from initial state $(a_0, c_0) = (2^k, r)$ until $a_t$ becomes odd (case iii).

\begin{lemma}[Prefix determinism]
\label{lem:prefix-determinism}
Let $n=2^k \cdot m + r$ with $r$ odd. The prefix iteration starting from initial state $(a_0, c_0) = (2^k, r)$ is fully determined by $r$: the sequence of state transitions, the terminal state $(a_{\text{final}}, c_{\text{final}})$, and the number of prefix steps are all independent of $m$.
\end{lemma}

\begin{proof}
The state transitions (i) and (ii) depend on $c_t \bmod 2$ and the parity of $a_t$. Both are functions of the iteration history starting from $(2^k, r)$. By the parity rules above, while $a_t$ is even the parity of $a_t \cdot m + c_t$ equals $c_t \bmod 2$ independent of $m$. The prefix iteration consists exactly of state transitions (i) and (ii), which apply while $a_t$ is even and terminate when $a_t$ becomes odd. Therefore the entire prefix iteration -- sequence of transitions, terminal state, and number of prefix steps -- is a function of $r$ alone.
\end{proof}

\subsection{The termination theorem}
\label{sec:termination}

\begin{theorem}[Termination]
\label{thm:termination}
The prefix iteration terminates in finitely many steps. The number of even-step applications (case i) consumed during the prefix is exactly $k$. The number of odd-step applications (case ii) is some $j$ with $1 \leq j \leq k$. The terminal $a_{\text{final}}$ is $3^j$.
\end{theorem}

\begin{proof}
The initial value $a_0 = 2^k$ has $v_2(a_0) = k$, where $v_2$ denotes the $2$-adic valuation. Even-step applications (case i) replace $a_t$ by $a_t / 2$ and decrement $v_2(a_t)$ by $1$. Odd-step applications (case ii) replace $a_t$ by $3 a_t$; since $3$ is odd, $v_2(a_t)$ is unchanged. Therefore $v_2(a_t)$ is monotone non-increasing along the prefix iteration. It starts at $k$ and decreases by exactly $1$ with each case-(i) step. The iteration enters case (iii) the first time $v_2(a_t) = 0$, i.e., the first time $a_t$ is odd. The number of case-(i) steps consumed during the prefix is therefore exactly $k$.

For the lower bound $j \geq 1$ follows from the initial state: $a_0 = 2^k$ is even and $c_0 = r$ is odd, so the first step is case (ii), giving $j \geq 1$.

For the upper bound $j \leq k$: after any case-(ii) step, $c_{t+1} = 3c_t +1$ is even (since $c_t$ is odd, $3 c_t$ is odd, and $3 c_t + 1$ is even), so the next step is necessarily case (i). Odd-steps cannot occur consecutively. The prefix consumes exactly $k$ case-(i) steps before termination (by the $2$-adic valuation argument). Each odd-step is immediately followed by an even-step, so the number of odd-steps is at most the number of even-steps, giving $j \leq k$.
    
The terminal $a_{\text{final}} = 3^j$, since each odd-step application multiplies $a$ by $3$ and the terminal value of $a$ (after the iteration's $k$ even-step applications have decremented the 2-adic valuation to $0$) is $3^j$ with $j$ the count of odd-step applications.
\end{proof}

The theorem is the deductive content of Lemma 1 plus the parity-rule arithmetic. Its proof is one paragraph and carries no heuristic content. The prefix iteration is well-defined, terminates in at most $2k$ steps total (at most $k$ even-step and at most $k$ odd-step applications), and produces a terminal $a_{\text{final}}$ of the form $3^j$ for $j \in \{1, 2, \ldots, k\}$.

A consequence of Theorem 1 is that the $2^{(k-1)}$ odd residue classes $\bmod 2^k$ partition into at most $k$ prefix-terminal types according to the value of $a_{\text{final}}$. Whether all $k$ types are actually realized at given $k$ (i.e., whether for each $j \in \{1, \ldots, k\}$ there exists $r \bmod 2^k$ whose prefix terminates at $a_{\text{final}} = 3^j$) is a separate question.

\subsection{Realization of all $k$ values of $a_{\text{final}}$}
\label{sec:realization}
\begin{corollary}[Realization]
\label{cor:realization}
For every $k \geq 1$ and every $j \in \{1, \ldots, k\}$ there exists an odd residue class $r \bmod 2^k$ with $a_{\text{final}}(r) = 3^j$.
\end{corollary}

\begin{proof}
    By Lemma~\ref{lem:j-binomial}, the number of such classes is $\binom{k-1}{j-1} \geq 1$.
\end{proof}

\subsection{The closed-form predictor \texorpdfstring{$\alpha_{\text{det}}$}{alpha-det}}
\label{sec:closed-form}
We define the per-class structural offset $\alpha_{\text{det}}(r)$ by
$$\alpha_{\text{det}}(r) := \text{prefix\_steps}(r) + K_h \cdot \log(a_{\text{final}}(r)/2^k) = (k + j(r)) + K_h \cdot (j(r) \log 3 - k \log 2),$$
where $j(r)$ is the number of odd-step applications consumed during prefix iteration starting from $r$, $\text{prefix\_steps}(r) = k + j(r)$, $a_{\text{final}}(r) = 3^{j(r)}$, and $K_h = 3/\log(4/3)$. The quantity $\alpha_{\text{det}}(r)$ is a closed-form function of $r$ alone: both $j(r)$ and $\text{prefix\_steps}(r)$ are deterministic outputs of the prefix iteration (Lemma~\ref{lem:prefix-determinism}, Theorem~\ref{thm:termination}).

The form of $\alpha_{\text{det}}$ is motivated by the standard random-walk heuristic \citep{lagarias1985}: Applying $K_h \log(\cdot)$ to the post-prefix state $a_{\text{final}} \cdot m + c_{\text{final}}$ recovers the leading-order behavior of $\sigma$ on the post-prefix orbit. This separates the deterministic contribution of the prefix (which depends only on $r$) from the stochastic contribution of the post-prefix dynamics (which depends on $m$). This is the only heuristic input to the construction; we present \S\ref{sec:empirics} as empirical validation of the per-class predictive content of this heuristic.

\begin{lemma}[Block decomposition]
\label{lem:block-decomposition}
The prefix trace of any odd $r \bmod 2^k$ decomposes uniquely as a sequence of $k$ blocks, each of one of two types. Type $O$ is a case-(ii) step immediately followed by a case-(i) step (the substring \texttt{oe}). Type $E$ is a single case-(i) step (the substring \texttt{e}). The first block of every prefix trace is of type O, and the number of type-O blocks equals $j$, the number of odd-step applications in the prefix.
\end{lemma}

\begin{proof}
    By the argument in the proof of Theorem~\ref{thm:termination}, after any case-(ii) step $c_{t+1} = 3 c_t + 1$ is even, so the next step is case (i): odd-steps cannot occur consecutively. Grouping each odd-step with its successor partitions the trace into type-O blocks ("oe") and type-E blocks ("e"). The first block is type O because $c_0 = r$ is odd. By Theorem~\ref{thm:termination} the prefix consumes exactly $k$ case-(i) steps in total; each block contains one case-(i) step, so the total number of blocks is $k$. The number of type-O blocks equals the number of odd-steps, $j$.
\end{proof}

\begin{lemma}[Block sequence equals Syracuse parity vector]
\label{lem:syracuse-equivalence}
Let $S \colon \mathbb{N} \to \mathbb{N}$ denote the Syracuse compressed map $S(c) = (3c+1)/2$ if $c$ is odd and $S(c) = c/2$ if $c$ is even. Define $c_0 := r$ and $c_{t+1} := S(c_t)$ for $t \geq 0$. Each block of the prefix trace at level $k$ corresponds to one application of $S$: the $t$-th block is type O if $c_t$ is odd and type E if $c_t$ is even. The block sequence of the prefix is the length-$k$ Syracuse parity vector $\Pi_k(r) := (c_0 \bmod 2, c_1 \bmod 2, \ldots, c_{k-1} \bmod 2)$.
\end{lemma}

\begin{lemma}[Binomial $j$-distribution]
\label{lem:j-binomial}
Among the $2^{k-1}$ odd residue classes $r \bmod 2^k$, the number whose prefix iteration terminates with exactly $j$ odd-step applications is $\binom{k-1}{j-1}$, for $j \in \{1, 2, \ldots, k\}$. Consequently $\langle j \rangle = (k+1)/2$.
\end{lemma}

\begin{proof}
    By Lemma~\ref{lem:syracuse-equivalence}, the block sequence of the prefix at level $k$ equals the Syracuse parity vector $\Pi_k(r)$, and the number of type-O blocks equals the Hamming weight of $\Pi_k(r)$. The classical Syracuse parity-vector bijection (\citet{terras1976}, Lemma 1; \citealt{lagarias1985}, \S 2) states that $\Pi_k \colon \mathbb{Z}/2^k\mathbb{Z} \to \{0,1\}^k$ is a bijection. Restricting to odd $r$ corresponds to fixing the first coordinate $c_0 \bmod 2 = 1$, giving a bijection between odd residues mod $2^k$ and $\{1\} \times \{0,1\}^{k-1}$. The number of vectors in this restricted set with Hamming weight $j$ (i.e., with $j-1$ ones among the $k-1$ free coordinates) is $\binom{k-1}{j-1}$.

    The mean is
    \begin{align*}
    \langle j \rangle = 2^{-(k-1)} \sum_{j=1}^{k} j \binom{k-1}{j-1} = 2^{-(k-1)} \sum_{i=0}^{k-1} (i+1) \binom{k-1}{i} = 2^{-(k-1)} \left[ (k-1) 2^{k-2} + 2^{k-1} \right] = (k+1)/2
    \end{align*}
\end{proof}

\begin{theorem}[Closed form for $\langle \alpha_{\text{det}} \rangle$]
\label{thm:alpha-det-closed-form}
Under the random-walk heuristic ($K_h = 3/\log(4/3)$) applied to the post-prefix state and Lemma~\ref{lem:j-binomial},
$$\langle \alpha_{\text{det}} \rangle = \frac{\log 6}{\log(4/3)}, \quad \text{independent of } k.$$
\end{theorem}

\begin{proof}
    From $\alpha_{\text{det}}(r) = \text{prefix\_steps}(r) + K_h \log(a_{\text{final}}(r)/2^k)$, Theorem~\ref{thm:termination}'s $\text{prefix\_steps}(r) = k + j(r)$:
    $$\alpha_{\text{det}}(r) = (k + j) + K_h (j \log 3 - k \log 2).$$
    Averaging over odd $r \bmod 2^k$ and using $\langle j \rangle = (k+1)/2$ from Lemma~\ref{lem:j-binomial}:
    $$\langle \alpha_{\text{det}} \rangle = k(1 - K_h \log 2) + \frac{k+1}{2}(1 + K_h \log 3).$$
    Writing $L = \log(4/3) = 2 \log 2 - \log 3$ and using $K_h = 3/L$:
    \begin{align*}
    1 - K_h \log 2 &= \frac{L - 3 \log 2}{L} = -\frac{\log 6}{L}, \\
    1 + K_h \log 3 &= \frac{L + 3 \log 3}{L} = \frac{2 \log 6}{L}.
    \end{align*}
    Therefore
    \begin{align*}
    \langle \alpha_{\text{det}} \rangle
      &= k \cdot \left(-\frac{\log 6}{L}\right) + \frac{k+1}{2} \cdot \frac{2 \log 6}{L} \\
      &= \frac{\log 6}{L}\bigl[-k + (k+1)\bigr] \\
      &= \frac{\log 6}{\log(4/3)}.
    \end{align*}
\end{proof}

Theorem~\ref{thm:alpha-det-closed-form} is exact in the prefix algebra and Lemma~\ref{lem:j-binomial}. The only non-deductive input is the application of $K_h$ from the random-walk argument to the post-prefix state. Direct enumeration of $\alpha_{\text{det}}(r)$ across all $2^{k-1}$ odd classes at $k \in \{6, 8, 10, 12, 14\}$ recovers $\log 6 / \log(4/3) \approx 6.228263$ to machine precision (relative error $\leq 10^{-14}$, as verified in \S\ref{sec:closed-form-verification}).

\section{Empirical predictions of the prefix decomposition}
\label{sec:empirics}
The findings in this section are empirical regularities verified at modular resolutions $k \in \{6, 7, 8, 9\}$ and $N$ up to $2^{27}$. Each finding is named, calibrated, and bounded by what was actually tested.

\subsection{Per-class intercept $\alpha(r)$}
At modular resolution $k = 6$ with $N=2^{25}$, hierarchical Bayesian fit of $\sigma \sim \alpha(r) + \beta(r) \cdot \log(n) + \varepsilon$ with class-level $\alpha$ and $\beta$ yields per-class $\alpha$ posterior means spanning $-27$ to $+35$ with $\tau_\alpha \approx 11.7$ (\texttt{writeup.md Result 3}). Linear regression of these posterior means against $\alpha_{\text{det}}(r)$ gives $\alpha_{\text{actual}}(r) = -2.66 + 0.987 \cdot \alpha_{\text{det}}(r)$, $R^2 = 0.9996$, with $\mathrm{SD}(\text{per-class residual}) = 0.28$ and $\mathrm{SD}(\text{residual})/\mathrm{mean}(\text{per-class SE}) = 0.18$. The residuals are below the noise floor of posterior sampling alone.

Extension to higher modular resolutions at $N=2^{27}$ via direct OLS regression (per-class regression rather than hierarchical Stan, to avoid posterior-geometry difficulties at large $k$):

\begin{table}[ht]
\centering
\caption{Per-class residual scale calibration across modular resolutions.}
\label{tab:residual-calibration}
\begin{tabular}{rrrrrc}
\toprule
$k$ & mod & classes & per class $n$ & $R^2$ & $\text{SD(resid)} / \text{mean(SE)}$ \\
\midrule
6 & 64 & 32 & 2.10M & 0.9967 & 0.96 \\
7 & 128 & 64 & 1.05M & 0.9942 & 0.99 \\
8 & 256 & 128 & 524K & 0.9918 & 0.91 \\
9 & 512 & 256 & 262K & 0.9851 & 0.93 \\
\bottomrule
\end{tabular}
\end{table}

\begin{observation}[Per-class intercept prediction]
\label{obs:intercept}
At modular resolutions $k \in \{6, 7, 8, 9\}$ and $N$ up to $2^{27}$, per-class intercept $\alpha(r)$ of $\sigma$ vs $\log(n)$ is predicted by $\alpha_{\text{det}}(r)$ with regression $R^2 \geq 0.985$ and signal-to-noise ratio $\text{SD(residual)}/ \text{mean(SE)}$ in $[0.91, 0.99]$. Per-class residuals scale with sampling noise, with no detectable signal above sampling noise at the cited resolutions. The regularity is verified at the cited resolutions; we do not claim it for arbitrary $k$ or for $k > 9$.

\end{observation}

\paragraph{Controlled comparison: matched per-class sample size.}
The decline in $R^2$ across $k \in \{6, 7, 8, 9\}$ in Table~\ref{tab:residual-calibration} reflects shrinking per-class sample sizes, not weakening of $\alpha_{\text{det}}$ as a predictor. At fixed per-class $n$, the $R^2$ is essentially independent of $k$:

\begin{center}
\begin{tabular}{rcrr}
\toprule
 $k$ & $N$ & per-class $n$ & $R^2$ \\
 \midrule
 6 & $2^{25}$ & 524K & 0.9967 \\
 8 & $2^{27}$ & 524K & 0.9918 \\
 \bottomrule
\end{tabular}
\end{center}

Same per-class sample size at different modular resolutions produces essentially the same $R^2$. The $k$-gradient in Table~\ref{tab:residual-calibration} is sampling noise, not signal degradation: as $k$ grows, the $2^{k-1}$ classes share the same total $N$, so per-class $n$ shrinks by half at each $k$, and per-class $\alpha$ posterior SE grows correspondingly. The signal-to-noise band $[0.91, 0.99]$ in Table~\ref{tab:residual-calibration} confirms residuals are at the noise floor across all four resolutions.

\subsection{The prefix-collapse onto $k$ distinct distributions}
\label{sec:moment-clustering}
The hypothesis initially tested was full distributional universality of $S(n) := \sigma(n) - \alpha(r) - \beta \cdot \log n$, i.e., that $S(n)$ would be class-independent. It is not. Per-class variance differs by approximately $30\%$ across classes mod $64$, and per-class skewness and kurtosis also vary. However, each higher central moment, when computed per class and plotted against $\alpha_{\text{det}}$ (or equivalently against $\log(a_{\text{final}})$, since these are collinear at fixed $k$), shows the prefix-prediction pattern:

\begin{table}[ht]
\centering
\caption{Pearson correlations of $\alpha_{\text{det}}$ with per-class central moments}
\label{tab:central-moments}
\begin{tabular}{rrrrrr}
\toprule
$k$ & classes & per-class $n$ & $r(\alpha_{\text{det}}, \text{Var})$ & $r(\alpha_{\text{det}}, \text{Kurt})$ & $r(\alpha_{\text{det}}, \text{Skew})$ \\
\midrule
6 & 32 & 2.10M & 0.99989 & 0.9497 & 0.8716 \\
7 & 64 & 1.05M & 0.99977 & 0.9337 & 0.8442 \\
8 & 128 & 524K & 0.99969 & 0.9236 & 0.8250 \\
9 & 256 & 262K & 0.99943 & 0.9077 & 0.7914 \\
\bottomrule
\end{tabular}
\end{table}

The clustering is sharper than the Pearson correlations alone indicate. When per-class variance is plotted against $\alpha_{\text{det}}$, classes with the same $a_{\text{final}}$ values produce 6 distinct clusters; at $k=7$, the 7 $a_{\text{final}}$ values produce 7 clusters; and so on through $k=9$. The same is true for kurtosis. Skewness shows slightly more within-cluster scatter but the same discrete structure.

\begin{observation}[Moment clustering by $a_{\text{final}}$]
\label{obs:moment-clustering}
At modular resolutions $k \in \{6, 7, 8, 9\}$, per-class central moments of $\sigma$ residuals cluster by $a_{\text{final}}$, with the number of distinct clusters equal to $k$ at every tested $k$. The regularity is verified at the cited resolutions.
\end{observation}

\subsection{Numerical verification of Theorem~\ref{thm:alpha-det-closed-form}}
\label{sec:closed-form-verification}
Theorem~\ref{thm:alpha-det-closed-form} predicts $\langle \alpha_{\text{det}} \rangle = \log 6 / \log(4/3) \approx 6.228263$ exactly, independent of $k$. Direct computation of $\alpha_{\text{det}}(r)$ across all $2^{k-1}$ odd residue classes (no sampling) recovers this value to machine precision at every tested resolution.

\begin{center}
\begin{tabular}{rrr}
\toprule
$k$ & $2^{k-1}$ classes & $\langle \alpha_{\text{det}} \rangle - \log 6 / \log(4/3)$  \\
\midrule
6 & 32 & $-5.3 \times 10^{-15}$ \\
8 & 128 & $-7.1 \times 10^{-15}$ \\
10 & 512 & $-7.1 \times 10^{-15}$ \\
12 & 2048 & $-8.9 \times 10^{-15}$ \\
14 & 8192 & $-1.1 \times 10^{-14}$ \\
\bottomrule
\end{tabular}
\end{center}

This is a sanity check on Theorem~\ref{thm:alpha-det-closed-form} rather than experimental evidence for it. The closed form is a deductive consequence of Lemma~\ref{lem:j-binomial} and the definition of $\alpha_{\text{det}}$, both of which are non-empirical inputs. We report the verification at $k \leq 14$ because beyond that the enumeration over $2^{k-1}$ classes becomes expensive, not because the result is range-bounded; the theorem holds for all $k$.

\subsection{Rate of convergence and limitations}
The experimental regularities are at finite $N$ (up to $2^{27}$ for the $k$-sweep). The $k$-sweep at fixed $N=2^{27}$ confirms the regularities at four distinct modular resolutions; we do not extrapolate to arbitrary $k$. Higher-$k$ extension is in principle possible but per-class sample sizes shrink as $2^{(1-k)}$; meaningful verification at $k > 12$ would require $N$ substantially larger than $2^{32}$.

The signal-to-noise ratio $\text{SD(residual)}/ \text{mean(SE)}$ for the $\alpha_{\text{det}}$ predictions is in the band $[0.91, 0.99]$ across $k \in \{6, 7, 8, 9\}$. At $k=4$ the ratio is $0.70$ -- residuals smaller than per-class SE, meaning the prefix prediction over-explains at coarse resolution (signal-to-noise below sampling noise). At every $k$ tested, the regularities of \S4.1 and \S4.2 hold within their stated bounds. We do not claim the regularities hold at modular resolutions outside the tested range.

\section{Relation to Terras (1976)}
\label{sec:terras-relation}
The prefix decomposition relates to \citet{terras1976} in two places.

First, the proof of Lemma~\ref{lem:j-binomial} relies on \citet{terras1976} Lemma 1, the Syracuse parity-vector bijection $\Pi_k \colon \mathbb{Z}/2^k\mathbb{Z} \to \{0,1\}^k$. The block decomposition of the prefix trace (Lemma~\ref{lem:block-decomposition}) and its identification with the Syracuse parity vector (Lemma~\ref{lem:syracuse-equivalence}) reduce the count of odd classes with $j$ odd-step applications to counting parity vectors of given Hamming weight under Terras's bijection. The prefix-algebra content of this paper rests on Terras 1976 Lemma 1.

Second, the prefix decomposition is a constructive sharpening of \citet{terras1976} Lemma 4. Terras's Lemma 4 gives the asymptotic ratio $S_k \approx S_0 \cdot 3^{d(k)}/2^k$ for the $k$th iterate's residue-class structure but does not retain the symbolic state past the asymptotic identity. The prefix decomposition retains the symbolic state $(a, c)$ through the deterministic prefix, terminates at $a_{\text{final}} = 3^j$ explicitly, and exposes $a_{\text{final}}$ and $\text{prefix\_steps}$ as closed-form covariates of per-class $\sigma$ statistics at finite $k$. The constructive content is what enables the per-class predictions of \S\ref{sec:empirics}; Terras's asymptotic content, by itself, does not produce them. We do not claim novelty relative to \citet{terras1976}; the prefix decomposition is a constructive use of Terras 1976 Lemma 1 to give a finite-$k$ refinement of Terras 1976 Lemma 4.

\section{Methods}

\subsection{Use of AI tools}
During the preparation of this work the author used Anthropic's Claude to assist with LaTeX coding, formatting and mathematical notation. The author wrote and edited the content and takes full responsibility for the content of the publication. The tool was not used to generate mathematical results, proofs, or conjectures, all of which are the author's own.

\section{Discussion}
\label{sec:discussion}
\subsection{What the empirical regularities suggest}
At every tested modular resolution $k \in \{6, 7, 8, 9\}$, the prefix decomposition's $\alpha_{\text{det}}$ predicts per-class statistics of $\sigma$ at the level of sampling noise. The signal-to-noise ratio $\text{SD(Residual)}/ \text{mean(SE)}$ is in $[0.91, 0.99]$ across four resolutions. Per-class moments cluster by $a_{\text{final}}$ at every tested $k$, with the number of distinct clusters equal to $k$.

\begin{conjecture}[Prefix parameterization of $\sigma \mid r$]
\label{conj:prefix-parameterization}
For every $k \geq 1$, the conditional distribution of the total stopping time $\sigma(n)$ given $r = n \bmod 2^k$ is parameterized by $a_{\text{final}}(r)$ alone among prefix-derivable quantities. Specifically:
\begin{enumerate}
    \item[(a)] (Mean.) $\mathbb{E}[\sigma(n) \mid r] = \alpha_{\text{det}}(r) + K_h \log n + o(1)$ as $n \to \infty$ within the class $r \bmod 2^k$.
    \item[(b)] (Higher moments.) For each $\ell \geq 2$, the $\ell$th central moment of $\sigma(n) \mid r$ depends on $r$ only through $a_{\text{final}}(r)$, asymptotically.
\end{enumerate}
\end{conjecture}

The experimental content of \S\ref{sec:empirics} supports both clauses at the tested resolutions. Clause (a) is the content of Observation~\ref{obs:intercept}: the per-class intercept is recovered by $\alpha_{\text{det}}$, with residuals at the noise floor of posterior sampling. Clause (b) is the content of Observation~\ref{obs:moment-clustering}: per-class higher moments cluster discretely by $a_{\text{final}}$ into exactly $k$ groups, with within-cluster variation at sampling-noise scale. A proof of either clause is open. We formulate the conjecture so that future work can confirm, refute, or sharpen it precisely.

\subsection{What is required for a stronger result}
Two pieces of analytical work would strengthen the empirical regularities of \S\ref{sec:empirics} toward theorem-grade results:
\begin{enumerate}
    \item Analytical characterization of the post-prefix dynamics from state $(a_{\text{final}}, c_{\text{final}}, m)$. The standard Collatz random-walk heuristic gives leading-order behavior on the post-prefix state but does not capture the higher-moment structure that the prefix-collapse predicts experimentally. A characterization that recovers the empirical Pearson correlations of \S\ref{sec:moment-clustering} (per-class variance, kurtosis, skewness clustering by $a_{\text{final}}$) from the post-prefix dynamics analytically would be a substantial strengthening.
    \item Connection to \citep{tao2022}'s almost-all-$N$ theorem. Preliminary experimental work in the accompanying repository extends $\alpha_{\text{det}}$ to predict per-class mean first-passage time at thresholds $f(N)$, of the form $s_{\text{mean}}(r; f) \approx \alpha_{\text{det}}(r) + K_h \log(N / f(N))$. This extension is reported separately and not claimed here. A theorem-grade version would derive the per-class structural offset $\alpha_{\text{det}}$ from Tao's framework directly, connecting the prefix decomposition to the leading-term of Tao's (5.15).
\end{enumerate}

None of these are claimed in the present paper.

\subsection{Extensions to $qx+1$}
The prefix decomposition formally extends to $qx+1$ dynamics on the residue class $\bmod 2^k$ for $q$ coprime to 2. The terminal $a_{\text{final}}$ lives in $\{ q^j : 1 \leq j \leq k \}$. The construction's algebraic content -- Lemma 1 (parity determined by $r$ alone), Theorem 1 (termination within at most $k$ odd-step applications), and the structure of the prefix iteration -- transfers directly with $q$ in place of $3$.

Whether the empirical regularities of \S4 transfer to $qx+1$ is an open question. Preliminary experimental work at $q \in \{5, 7, 9, 11\}$ is reported separately. The $3x+1$ results of the present paper are restricted to the $q=3$ case.

\section{Conclusion}
\label{sec:conclusion}
This paper has established three classes of result on the prefix decomposition of the Collatz map.
 
\subsection{Algebraic content} The prefix iteration is well-defined (Lemma~\ref{lem:prefix-determinism}), terminates with exactly $k$ even-step applications and at most $k$ odd-step applications (Theorem~\ref{thm:termination}), and produces a terminal $a_{\text{final}} \in \{3^j : 1 \leq j \leq k\}$. The $j$-distribution across the $2^{k-1}$ odd classes mod $2^k$ is exactly binomial: $\binom{k-1}{j-1}$ classes terminate at $a_{\text{final}} = 3^j$ (Lemma~\ref{lem:j-binomial}), with realization of every $j \in \{1, \ldots, k\}$ as a corollary (Corollary~\ref{cor:realization}). Theorem~\ref{thm:alpha-det-closed-form} establishes $\langle \alpha_{\text{det}} \rangle = \log 6 / \log(4/3)$ exactly in the prefix algebra, with the random-walk constant $K_h$ as the only heuristic input.
 
\subsection{Experimental content} At modular resolutions $k \in \{6, 7, 8, 9\}$ and $N$ up to $2^{27}$, the random-walk embedding $K_h = 3/\log(4/3)$ underlying $\alpha_{\text{det}}$ is validated at the per-class level. The per-class intercept $\alpha(r)$ of $\sigma$ against $\log n$ is predicted by $\alpha_{\text{det}}(r)$ with $R^2 \geq 0.985$ and residual signal-to-noise ratio $\mathrm{SD}(\text{residual})/\mathrm{mean}(\mathrm{SE}) \in [0.91, 0.99]$. Per-class central moments of $\sigma$ residuals cluster by $a_{\text{final}}$ into exactly $k$ distinct groups at every tested $k$. These regularities are verified at the cited resolutions and are not extrapolated.
 
\subsection{Conjectural content} The experimental regularities motivate Conjecture~\ref{conj:prefix-parameterization}: that the conditional distribution $\sigma \mid r \bmod 2^k$ is parameterized by $a_{\text{final}}$ alone at arbitrary $k$, both in its mean and in its higher moments. A proof of either clause is open.
 
The prefix-algebra results of \S\ref{sec:prefix} constitute a constructive sharpening of \citet{terras1976} Lemma 4 using Terras's own Lemma 1 (the Syracuse parity-vector bijection).

Two open questions remain: analytical characterization of the post-prefix conditional distribution that recovers the empirical clustering of per-class central moments by $a_{\text{final}}$, and the formal connection of $\alpha_{\text{det}}$ to \citet{tao2022}'s almost-all-$N$ theorem at the level of per-class structural offset.

\section*{Statements and Declarations}

\subsection*{Funding}
The author received no funding for this work.

\subsection*{Competing Interests}
The author has no competing interests to declare that are relevant to the
content of this article.

\subsection*{Data Availability}
All code, data, and computational records supporting the results and analysis
in this article are openly available in the accompanying repository at
\url{https://github.com/mrnathanhumphrey-droid/Collatz} and archived on Zenodo
at \url{https://doi.org/10.5281/zenodo.20100517}. The specific scripts
referenced in the text are listed in Appendix~\ref{app-reprod}.

\appendix

\section{Algorithmic specification of the prefix iteration}
\label{app:prefix-algorithm}
The prefix iteration is implemented directly from the parity rules \S\ref{subsec:prefix-recall}. Pseudocode for computing $(a_{\text{final}}, c_{\text{final}}, \text{prefix\_steps})$ given odd residue $r$ and modular resolution $k$:

\begin{verbatim}
function prefix(r, k):
    a := 2^k
    c := r
    steps := 0
    while a is even:
        if c is even:
            a := a / 2
            c := c / 2
        else:
            a := 3 * a
            c := 3 * c + 1
        steps := steps + 1
    return (a, c, steps)
\end{verbatim}

The iteration terminates when $a$ becomes odd (the loop's exit condition). Theorem~1 guarantees termination within at most $2k$ iterations. Reference implementations in the accompanying repository: \texttt{alpha\_beta\_gamma\_decay.py} and \texttt{analytical\_abc\_derivation.py}.

\section{Reproducibility}\label{app-reprod}
\label{app:reproducibility}
All code, data, and computational records are available at the accompanying repository \citep{humphrey_collatz_2026}. Specific scripts referenced in this paper:

\begin{itemize}
    \item experiments/01\_alpha\_decomposition.py -- $k=6$ hierarchical Bayesian fit at $N = 2^{25}$.
    \item experiments/24\_k\_sweep\_alpha\_decomposition.py -- $k$-sweep at $N = 2^{27}$ for $k \in \{4, \ldots, 12\}$.
    \item experiments/09\_multi\_stat\_decomposition.py -- per-class higher central moments (variance, kurtosis, skewness).
\end{itemize}

\subsection{Conventions}
The total stopping time $\sigma(n)$ is defined as the number of Collatz steps required to reach $n = 1$ (not the steps-to-first-descent-below-$n$ variant). Iteration uses the unaccelerated Collatz map $T(n) = 3n+1$ if $n$ odd, $n/2$ if $n$ even, with no Syracuse compression. All computations of $\sigma$ and $\alpha_{\text{det}}$ are deterministic; no random sampling is used. Per-class $\alpha(r)$ is estimated by OLS regression of $\sigma(n)$ on $\log(n)$ within each residue class $r \bmod 2^k$ over $n \leq N$.

Detailed experimental records are in \texttt{findings.md} and \texttt{closed\_form\_findings.md} (the latter contains the derivation of Theorem~\ref{thm:alpha-det-closed-form} as Result 1) within the repository.




\bibliography{references}

\end{document}
